\documentclass[12pt,a4paper]{article}
\usepackage[utf8]{inputenc}\usepackage[T1]{fontenc}\usepackage[provide=*,french]{babel}
\usepackage{lmodern,microtype,longtable,booktabs,array,enumitem,fancyhdr,needspace}
\usepackage[margin=2.1cm,headheight=16pt]{geometry}
\usepackage[hidelinks,breaklinks=true]{hyperref}
\setlength{\parindent}{0pt}\setlength{\parskip}{6pt}\setlength{\emergencystretch}{3em}
\setlist[itemize]{leftmargin=1.3em,itemsep=3pt}
\pagestyle{fancy}\fancyhf{}\fancyhead[L]{\small Master Finance verte — Module 6}\fancyhead[R]{\small Badr Aboufarasse}\fancyfoot[C]{\thepage}
\begin{document}
\begin{center}{\small MASTER FINANCE VERTE -- MODULE 6}\par\vspace{8pt}{\Large\bfseries Séance 6 -- Décision intégrée et pilotage financier de la transition}\par\vspace{6pt}Badr Aboufarasse\par{\small Cours et applications corrigées -- 25 octobre 2026}\end{center}
\begin{quote}\itshape Quelle trajectoire financer, sous quelles conditions et avec quel suivi ?\end{quote}
\textbf{Données du cas fictives.} Montants en M MAD ; taux et seuils pédagogiques.\par Durée : 3 h 30 par séance ; 21 heures pour le module.

\subsection*{Objectifs}
\begin{itemize}\item Assembler une recommandation stratégique cohérente et traçable.\item Connecter engagements, investissements, financements et résultats physiques.\item Organiser le suivi des écarts et la révision des décisions.\end{itemize}
\subsection*{Déroulé de 3 h 30}
Problème initial : 15 min ; concepts : 55 min ; exemples résolus : 30 min ; applications : 65 min ; cas : 35 min ; synthèse : 10 min. Total : 210 min.

\section*{1. Revenir à la question initiale}
Le comité doit choisir une trajectoire dont les conditions de réussite sont identifiées. Les cinq séances précédentes ont établi le diagnostic, le calendrier des besoins, les offres de financement, les modalités de développement et les limites du plan face aux chocs. La recommandation finale doit montrer comment ces éléments se combinent.


Une trajectoire peut créer de la valeur dans le scénario central mais être trop fragile au regard des ressources disponibles. Une autre peut mieux protéger la liquidité mais exposer l’entreprise à une dépendance contractuelle. La décision ne se réduit pas au plus grand chiffre de VAN ; elle doit définir quels risques sont acceptés, transférés ou réduits. [BD]


Dans le cas, il faut choisir un programme, un mode solaire et un financement, puis un dispositif de réaction. Toute alternative au programme de référence appelle une nouvelle prévision. Les tableaux de référence restent le point de comparaison et ne deviennent pas automatiquement les résultats de l’alternative retenue.

\section*{2. Construire une trajectoire de transition financée}
Un plan crédible relie objectifs, leviers opérationnels, ressources et suivi. Les ressources TPT organisent les informations autour des fondements, de la stratégie de mise en œuvre, de l’engagement, des métriques et objectifs, et de la gouvernance. Nous utilisons cette structure comme aide à la décision, sans la confondre avec une obligation générale de publication. [TPT]



\Needspace{180pt}
\begin{longtable}{p{3.67cm}p{3.67cm}p{3.67cm}p{3.67cm}}
\toprule
Objectif du cas & Levier & Ressource / calendrier & Preuve attendue \\
\midrule\endhead
Réduire l’énergie nécessaire & Efficacité E & 6 en A1 ; service A2 & Mesure avant/après ajustée à la production \\
Modifier l’approvisionnement & Solaire détenu ou contrat & Selon option et financement & Production, factures, attributs et disponibilité \\
Réduire la dépendance à l’eau & Projet W & 4 en A2 ; service A3 & Volumes, qualité et interruptions \\
\bottomrule
\end{longtable}

Le budget doit distinguer dépenses engagées et envisagées, financement confirmé et recherché. Un objectif à dix ans ne prouve pas que les cinq premières années sont financées. Les informations sur effets financiers et incertitudes doivent être cohérentes avec les prévisions. [IFRS]

\section*{3. Tableau de bord : résultats et signaux avancés}
Un bon tableau de bord est sélectif. Chaque indicateur précise définition, unité, périmètre, fréquence, source, responsable et seuil d’action. On distingue résultats constatés et signaux avancés : économie annuelle d’énergie d’un côté, retard d’installation ou disponibilité des compteurs de l’autre.



\Needspace{250pt}
\begin{longtable}{p{5.07cm}p{5.07cm}p{5.07cm}}
\toprule
Indicateur & Fréquence / responsable & Seuil pédagogique et action \\
\midrule\endhead
Cash minimum prévu sur 13 semaines & Hebdomadaire / trésorerie & < 3 : plan de liquidité \\
DSCR prévisionnel selon convention & Trimestrielle / DAF & < 1,20 : revue du financement \\
Mise en service E/S/W & Mensuelle / projets & Retard > 1 mois : nouveau calendrier de gains \\
Énergie par unité produite & Mensuelle / exploitation & Écart > 10 \% de la cible : analyse technique \\
Eau par unité et interruptions & Mensuelle / exploitation & Seuil local défini avant décision \\
\bottomrule
\end{longtable}

Les seuils 3, 1,20, un mois et 10 \% sont des choix du cas. Ils doivent être proportionnés et révisés au regard des engagements réels. Un indicateur sans source fiable peut créer une illusion de contrôle.

\section*{4. Intensité, absolu et périmètre environnemental}
Une intensité peut s’améliorer tandis que l’impact total augmente. Si la production passe de 100 à 120 et l’énergie par unité de 10 à 9, l’intensité diminue de 10 \% mais la consommation totale passe de 1 000 à 1 080, soit +8 \%. Suivre les deux valeurs évite de confondre efficacité et réduction absolue.


Pour le solaire du cas, nous ajoutons une hypothèse physique autonome : 3 GWh/an d’électricité de réseau substitués et un facteur pédagogique de 0,6 tCO2e/MWh. La quantité indicative est 3 000 $\times$ 0,6 = \textbf{1 800 tCO2e/an}. Ce facteur n’est pas une estimation officielle du réseau marocain. Le calcul ne couvre ni cycle de vie, ni scope 3, ni certification d’émissions évitées.


La comptabilité des émissions doit respecter les périmètres et les instruments applicables. Des attributs contractuels doivent être suivis pour éviter des allégations incompatibles ou des doubles revendications. On ne doit pas additionner mécaniquement une baisse d’inventaire et des « émissions évitées » comme deux réductions distinctes. [GHG]

\section*{5. Prix interne du carbone et arbitrage}
Un prix interne peut aider à comparer des investissements ou préparer une exposition future. Il faut préciser sa fonction : prix fictif d’aide à la décision, coût réglementaire prévu, ou prélèvement interne réellement transféré entre unités. Ces trois objets ne produisent pas les mêmes flux. [SS6]


Avec 1 800 tCO2e/an et un prix fictif de 500 MAD/t, la valorisation indicative atteint 0,9 M MAD/an. Elle n’entre pas dans le CFO en l’absence d’un mécanisme réel de paiement ou de recette. Si un coût carbone futur est déjà intégré à la facture évitée, le rajouter comme bénéfice financier produirait un double compte.


Le comité peut regarder une VAN financière et un résultat physique, puis appliquer une contrainte d’émissions ou une valeur d’impact séparée. Le cadre de valeur intégrée invite à examiner ces arbitrages ; il n’élimine pas la nécessité de justifier les prix, les pondérations et les effets distributifs. [SS2]

\section*{6. Gouverner les écarts et les décisions}
Le budget approuvé est une référence explicite, pas une interdiction d’apprendre. Lorsqu’un écart apparaît, le processus doit établir le fait, diagnostiquer la cause, recalculer les conséquences puis décider. La direction financière vérifie les flux et le financement ; l’exploitation valide les performances ; le responsable de projet met à jour les dates ; le comité approuve les modifications importantes.


Le lien avec le module 5 est opérationnel : objectif $\rightarrow$ mécanisme $\rightarrow$ contrôle $\rightarrow$ preuve $\rightarrow$ revue $\rightarrow$ action $\rightarrow$ nouveau test. Les responsabilités de production de l’information et de vérification doivent être identifiées. Les coûts d’agence éclairent l’intérêt des engagements et des contrôles, sans supposer qu’un reporting volumineux garantit la bonne décision. [JM]


Une décision conditionnelle peut autoriser E et S sous réserve d’un contrat confirmé, d’un échéancier de paiement et d’une réserve de liquidité ; elle peut prévoir une revue de W avant engagement si la contrainte hydrique le permet. Elle doit mentionner ce qui ferait abandonner ou accélérer l’option, et qui peut l’autoriser.

\section*{7. Approfondissement --- Distinguer budget initial, budget flexible et nouvelle prévision}
Le budget initial constitue la référence approuvée. Le budget flexible recalcule certains postes au niveau d’activité effectivement réalisé, pour distinguer l’effet du volume des autres écarts. La nouvelle prévision projette l’avenir avec les informations disponibles ; elle n’efface pas la référence initiale. Les trois objets répondent à des questions différentes et doivent être nommés.


Une économie solaire inférieure au budget peut venir de la production électrique, du prix évité, de la disponibilité ou de la maintenance. La décomposition doit respecter une convention pour l’interaction prix-volume. En appliquant d’abord le volume puis le prix sur le volume réalisé, on obtient une somme exacte des effets. Une autre convention est possible, à condition de ne pas attribuer deux fois l’interaction.


Les mêmes exigences valent pour les indicateurs physiques. Comparer l’intensité énergétique sans contrôler la production, le mix et la saison peut attribuer à un projet un gain qui vient en réalité d’un changement d’activité. Inversement, une hausse du total peut masquer une amélioration technique lorsque la production augmente. Suivre l’absolu, l’intensité et les facteurs explicatifs. [GHG]


Une information robuste doit préciser origine, date, périmètre, unité et retraitements. La prévision révisée doit ensuite répercuter les écarts sur le compte de résultat, la trésorerie, les engagements et le calendrier de décision. Le reporting devient utile lorsqu’il permet cette action, plutôt que la seule présentation d’une cible.

\section*{8. Approfondissement --- Distribution, réserve et décision multicritère}
Le cash disponible peut être affecté à l’investissement, à une réserve, au désendettement ou à une distribution. La décision dépend des perspectives, des risques et des engagements, pas uniquement du résultat net. Les travaux sur émissions et distributions rappellent cette concurrence entre usages des ressources et les tensions possibles avec les objectifs de transition. [SS16]


Un excédent de clôture du scénario central ne constitue pas automatiquement un montant distribuable. Il faut rapprocher échéances, contraintes contractuelles, besoins futurs et scénarios défavorables. Le calcul d’un maximum financier d’exercice ne vaut pas autorisation juridique de dividende ; les règles applicables et les décisions sociales ne sont pas analysées dans notre cas fictif.


Pour une décision multicritère, séparer \textbf{conditions d’admissibilité} et \textbf{préférences entre options admissibles}. Une bonne note environnementale ne permet pas de compenser arbitrairement un financement impossible. Après filtrage des contraintes, les pondérations doivent être explicites et leur sensibilité discutée. Un score agrégé facilite la comparaison mais ne remplace pas l’analyse des risques ni des données. [SS6]


\textbf{Structure d’une décision robuste.} Préciser le choix, les hypothèses décisives, le financement confirmé, les risques résiduels, les résultats attendus, les seuils de revue et la personne autorisée à modifier le programme. La règle doit aussi prévoir l’escalade lorsque les données nécessaires ne sont pas disponibles. Le comité approuve une trajectoire avec des conditions vérifiables.

\section*{9. Trois applications corrigées}

\subsection*{Application 1 --- Décomposer une économie manquante}
Le solaire devait éviter une dépense brute de 2,4. La production effective est 90 \% du prévu et le prix effectivement évité 95 \% du prévu. La maintenance reste 0,4. Comparer le gain réel au gain EBITDA prévu de 2.


\paragraph{Correction expliquée}
Gain brut réel = 2,4 $\times$ 0,90 $\times$ 0,95 = \textbf{2,052}. Gain net = \textbf{1,652} ; écart = -0,348. Décomposition séquentielle : volume -0,24 ; prix sur volume réel -0,108. La convention de décomposition doit rester constante pour attribuer l’interaction.



\subsection*{Application 2 --- Une incitation financière limitée}
Une remise de marge de 0,10 point de pourcentage s’applique pendant un an à un encours de 14 si l’objectif est atteint. La vérification coûte 0,03 pour cette même année. Ignorer fiscalité et actualisation. Quel avantage net ?


\paragraph{Correction expliquée}
Économie = 14 $\times$ 0,001 = \textbf{0,014} ; avantage net = \textbf{-0,016} M MAD. La remise seule ne couvre pas ce coût. D’autres bénéfices éventuels doivent être identifiés et ne peuvent être présumés.



\subsection*{Application 3 --- Écart absolu et intensité}
La production passe de 100 à 120 ; les émissions par unité de 2 à 1,8. Peut-on annoncer une baisse absolue de 10 \% ?


\paragraph{Correction expliquée}
Non. L’intensité baisse de 10 \% ; les émissions totales passent de 200 à \textbf{216}, soit +8 \%. L’annonce doit préciser indicateur, périmètre et année de référence.

\section*{10. Applications supplémentaires --- Écarts, objectifs et réserve de cash}

\subsection*{Application 4 --- Décomposer exactement un écart solaire}
Variante de performance : production solaire prévue 3 GWh, prix évité 0,80 MAD/kWh, maintenance 0,4 M MAD. Réalisé : 2,7 GWh, 0,85 MAD/kWh et maintenance 0,45. Calculer le gain net, l’écart total et la décomposition volume, prix et maintenance. Le produit de GWh par MAD/kWh donne ici directement des M MAD.


\paragraph{Corrigé détaillé : calcul, interprétation et limites}
Prévu : 3 $\times$ 0,80 - 0,4 = \textbf{2}. Réalisé : 2,7 $\times$ 0,85 - 0,45 = \textbf{1,845}. Écart = \textbf{-0,155}.

Effet volume au prix prévu : (2,7 - 3) $\times$ 0,80 = -0,24. Effet prix sur volume réalisé : 2,7 $\times$ (0,85 - 0,80) = +0,135. Effet maintenance : -(0,45 - 0,4) = -0,05. Total = -0,24 + 0,135 - 0,05 = -0,155.

Budget flexible au volume réalisé : 2,7 $\times$ 0,80 - 0,4 = \textbf{1,76}. L’écart réalisé/flexible est +0,085, tandis que l’effet du moindre volume est -0,24.

\textbf{Interprétation et limites.} Une hausse du prix évité peut masquer une production technique insuffisante. Le responsable d’exploitation doit analyser la production ; le DAF doit répercuter le gain réel et son profil temporel dans la prévision.


\subsection*{Application 5 --- Transformer une cible absolue en exigence d’intensité}
Consommation de référence 10 GWh pour une production de 100 unités. La production future atteint 125 unités. Une amélioration d’intensité de 25 \% est-elle suffisante pour atteindre une cible absolue de -20 \% ? Calculer la consommation, puis la baisse d’intensité nécessaire.


\paragraph{Corrigé détaillé : calcul, interprétation et limites}
Intensité initiale = 10/100 = 0,1 GWh/unité. Après amélioration de 25 \%, intensité = 0,075. Consommation future = 125 $\times$ 0,075 = \textbf{9,375 GWh}, soit seulement -6,25 \% en absolu.

La cible absolue est 8 GWh. Intensité requise = 8/125 = \textbf{0,064 GWh/unité}. Réduction nécessaire par rapport à 0,1 = \textbf{36 \%}. La baisse d’intensité de 25 \% est donc insuffisante de 11 points de pourcentage.

\textbf{Interprétation et limites.} Les unités de production et le périmètre doivent rester comparables. La consommation d’énergie et les émissions ne sont pas le même indicateur : un changement d’approvisionnement peut modifier le facteur d’émission sans la même variation de consommation.


\subsection*{Application 6 --- Combien de cash peut-on réellement distribuer ?}
La référence A1 se termine avec 4,72 de cash et une réserve souhaitée de 3. Le scénario défavorable atteint 1,834 en A1 puis 0,958 en A2. Aucun financement contingent n’est confirmé. Quel dividende maximal semble possible sur le seul A1 central ? Pourquoi cette conclusion change-t-elle si l’on exige de protéger la réserve dans le stress ?


\paragraph{Corrigé détaillé : calcul, interprétation et limites}
Sur la seule clôture centrale A1, plafond arithmétique = 4,72 - 3 = \textbf{1,72}. Mais le stress viole déjà la réserve sans distribution. En A2, l’écart à combler est 3 - 0,958 $\approx$ \textbf{2,042}, avant coûts d’un éventuel financement.

Une distribution immédiate aggraverait cette fragilité. Sous l’exigence de réserve au stress et sans autre action, aucun dividende positif n’est justifié par ces seuls chiffres ; une distribution nulle n’est même pas suffisante pour restaurer la réserve. Il faut modifier le plan ou confirmer des ressources.

\textbf{Interprétation et limites.} Le plafond central est une borne financière d’exercice, pas une autorisation légale ou contractuelle. Ne pas additionner un déficit de stress et un dividende comme s’ils appartenaient nécessairement à la même trajectoire réalisée ; il faut recalculer la politique retenue.


\subsection*{Application 7 --- Filtrer les contraintes avant de pondérer les préférences}
Trois options ont des notes pédagogiques sur 5, toutes orientées « plus est préférable » : détention (coût 4, liquidité 2, environnement 5, flexibilité 2), contrat (3, 5, 4, 3), phasage (3, 4, 3, 5). Pondérations 40 \%, 30 \%, 20 \%, 10 \%. Leurs DSCR sont respectivement 0,95 ; 1,30 ; 1,15, avec minimum obligatoire 1,20. Calculer les scores et l’admissibilité. Puis tester des poids 15 \%, 5 \%, 75 \%, 5 \%.


\paragraph{Corrigé détaillé : calcul, interprétation et limites}
Scores initiaux : détention = 1,6 + 0,6 + 1 + 0,2 = \textbf{3,4} ; contrat = 1,2 + 1,5 + 0,8 + 0,3 = \textbf{3,8} ; phasage = 1,2 + 1,2 + 0,6 + 0,5 = \textbf{3,5}.

Avec le minimum de DSCR, \textbf{seul le contrat est admissible}. Les nouveaux poids donnent détention 4,55 ; contrat 3,85 ; phasage 3,15. La détention devient première par préférence, mais reste inadmissible si sa structure financière n’est pas modifiée.

\textbf{Interprétation et limites.} Les notes, poids et DSCR sont des données d’une variante décisionnelle, pas des sorties du modèle de référence. Une pondération ne doit pas masquer une contrainte. Si l’on change le financement, recalculer le DSCR et les autres critères avant de reclasser les options.

\section*{11. Comité de synthèse : une décision exécutable}
\textbf{Atelier de clôture, 35 minutes, sans notation ni soutenance.} Une nouvelle information annonce un retard solaire et une hausse des coûts énergétiques. Rédiger une décision de dix lignes : programme retenu ; financement ; critères de valeur ; limite de liquidité ; résultats physiques ; conditions ; seuils ; responsables ; date de revue ; motif d’arrêt ou de révision.



\paragraph{Proposition de structure de réponse}
« Retenir E et le mode solaire choisi après validation de son coût complet et des engagements. Actualiser le plan avec le retard observé. Confirmer les financements avant acomptes et préserver un minimum de cash de 3. Réexaminer le scénario défavorable et la définition contractuelle du DSCR. Maintenir W selon la contrainte hydrique documentée. Confier au DAF une prévision glissante, aux opérations les mesures physiques et au comité les changements majeurs. Revoir les décisions aux seuils approuvés. » Cette structure doit être complétée par les chiffres de l’option effectivement recalculée ; elle ne démontre pas sa faisabilité à elle seule.


\textbf{Synthèse du parcours :} comprendre les contraintes ; dater les flux ; choisir les ressources ; comparer les engagements ; tester les chocs ; gouverner l’adaptation. Une bonne décision financière rend visibles les hypothèses, les tensions et les conditions qui pourraient changer la recommandation.


\textbf{Prolongements de lecture :} suivre les chapitres indiqués dans les références et comparer un plan de transition public à ses comptes et à son budget d’investissement. L’évaluation formelle du module sera construite ultérieurement.

\section*{Références mobilisées}
\small
\begin{itemize}
\item \textbf{[BD] Berk, J. et DeMarzo, P..} Corporate Finance, 6e édition, Pearson. Chapitres 2, 12, 14--18, 22, 24--27 et 30--31 selon la séance. \href{https://www.pearson.com/en-us/subject-catalog/p/corporate-finance/P200000009791/9780137844906}{Source officielle}.
\item \textbf{[TPT] Transition Plan Taskforce / IFRS Foundation.} Disclosure Framework et ressources sectorielles. \href{https://www.ifrs.org/sustainability/knowledge-hub/transition-plan-taskforce-resources/}{Source officielle}.
\item \textbf{[IFRS] IFRS Foundation (2025).} Disclosing information about an entity’s climate-related transition, including information about transition plans, in accordance with IFRS S2. \href{https://www.ifrs.org/content/dam/ifrs/supporting-implementation/ifrs-s2/transition-plan-disclosure-s2.pdf}{Source officielle}.
\item \textbf{[GHG] WRI / WBCSD (2015).} GHG Protocol Scope 2 Guidance, chap. 4 et critères de qualité des instruments contractuels. \href{https://ghgprotocol.org/sites/default/files/2023-03/Scope\%202\%20Guidance.pdf}{Source officielle}.
\item \textbf{[SS6] Schoenmaker, D. et Schramade, W. (2023).} Investment Decision Rules. Chap. 6, p. 145--171. \href{https://doi.org/10.1007/978-3-031-35009-2_6}{Source officielle}.
\item \textbf{[SS2] Schoenmaker, D. et Schramade, W. (2023).} Integrated Value Creation. Corporate Finance for Long-Term Value, chap. 2, p. 33--63. \href{https://doi.org/10.1007/978-3-031-35009-2_2}{Source officielle}.
\item \textbf{[JM] Jensen, M. C. et Meckling, W. H. (1976).} Theory of the firm: Managerial behavior, agency costs and ownership structure. Journal of Financial Economics, 3(4), 305--360. \href{https://doi.org/10.1016/0304-405X(76)90026-X}{Source officielle}.
\item \textbf{[SS16] Schoenmaker, D. et Schramade, W. (2023).} Issues and Payouts: Changes in Capital Structure. Chap. 16, p. 463--495. \href{https://doi.org/10.1007/978-3-031-35009-2_16}{Source officielle}.
\end{itemize}\normalsize
\end{document}
